\clearpage
\section[Rev. 1.42 -- Watchdog U403: disattivazione solo hardware]{Rev. 1.42 \textendash{} Watchdog U403: disattivazione solo hardware con R413, nessun GPIO su EN}

\subsection{Data e revisione}
03/11/2026, revisione 1.42.

\subsection{Titolo}
Watchdog U403: disattivazione solo hardware con R413, nessun GPIO su EN.

\subsection{Componenti, net e fogli coinvolti}
\begin{itemize}
\item Foglio MCU: U403 \code{STWD100NYWY3F}, R413 \SI{0}{\ohm} 0402 DNP (EN verso \code{3V3}).
\item PC6, rete \code{/MCU/WDT\_WDI}.
\item Nessun file di progetto modificato.
\end{itemize}

\subsection{Cosa è cambiato}
Nulla sullo schema. Decisione di progetto e precisazione della nota firmware F-11.

{\small\arrayrulecolor{rulec}
\begin{tabular}{@{}lll@{}}
\toprule
\textbf{Elemento} & \textbf{Prima (rev.\ 1.41)} & \textbf{Dopo (rev.\ 1.42)}\\
\midrule
Schema MCU & invariato & invariato\\
Disattivazione watchdog & R413 \SI{0}{\ohm} DNP & solo R413 (scelta A confermata)\\
GPIO su EN di U403 & non presente & non aggiunto (alternativa B scartata)\\
Nota firmware F-11 & impulso su WDI & precisata la modalit\`a di generazione dell'impulso\\
\bottomrule
\end{tabular}}

\subsection{Motivazione}
Brik81 ha verificato in KiCad il collegamento della rev.\ 1.41 (\textquotedblleft{}credo tutto ok\textquotedblright{}). Ha chiesto se collegare EN a un GPIO per disattivare il watchdog via software. Domanda testuale: \textquotedblleft{}se devo debuggare quidni questa resistenza va saldatata giusto? se si, forse sarebbe carino collegarlo al micro cosi posso disattivarlo via SW?\textquotedblright{}

Valutazione: un GPIO non sostituisce R413.
\begin{itemize}
\item Quando il firmware non gira (programmazione SWD di un chip vuoto, bootloader UART con BOOT0, firmware assente o rotto) il pin \`e flottante. Il pull-down interno di EN (\SIrange{32}{100}{\kilo\ohm}, datasheet ST DocID14134 Rev 10 p.\,12) tiene il watchdog attivo e l'MCU si resetta ogni circa \SI{1,6}{\second} durante la programmazione.
\item In volo un firmware fuori controllo potrebbe disattivare il proprio watchdog. Mitigabile con il lock GPIO (LCKR), non verificato sul reference manual dell'STM32L562.
\end{itemize}

Alternativa B scartata: GPIO libero (es.\ PC7) su EN tramite \SI{0}{\ohm} 0402 DNP. Scelta A: solo R413.
\begin{itemize}
\item Debug: si salda R413 (EN alto = watchdog spento). Si pu\`o tenere una scheda da banco con R413 sempre montata.
\item Volo: R413 non montata.
\end{itemize}

Approvazione di Brik81 (03/11/2026): \textquotedblleft{}No lascia comè e direi che siamo a posto giusto?\textquotedblright{}

\subsection{Impatti e criticit\`a}
\begin{notebox}
\textbf{Nessun impatto hardware.} Durante lo sviluppo R413 va montata: altrimenti debug SWD con breakpoint, bootloader UART e firmware di prova senza impulso su PC6 provocano reset ogni circa \SI{1,6}{\second}.
\end{notebox}

\begin{notebox}
\textbf{Nota firmware F-11 (precisazione).} L'impulso su WDI (PC6, \code{/MCU/WDT\_WDI}) va generato dal software nel ciclo principale, commutando il GPIO dopo aver verificato che il sistema funziona. NON va generato con un timer hardware in PWM o con un interrupt a tempo autonomo: il timer continua a girare anche con la CPU bloccata e il watchdog non scatterebbe mai.
\end{notebox}

\subsection{Verifiche da fare}
\begin{itemize}
\item M-04 invariata: con R413 DNP e firmware bloccato, impulso su NRST di circa \SI{210}{\milli\second} ogni circa \SI{1,6}{\second}; con R413 montata nessun reset.
\item In revisione del firmware controllare che F-11 sia rispettata (nessun PWM su PC6).
\end{itemize}
